\chapter[CONCLUSION AND FUTURE SCOPE]{Conclusion and Future Scope}

This chapter brings together the project outcomes and evidence mentioned in the report. It provides a summary of the contribution of the application, pinpoints the scope of the existing assessment, and defines the subsequent development of the data, explanation validation, and operations.

\section{Conclusions}
TrustFin AI links two product-specific loan classifiers with local explanations and diagnostic checks and brings it together to a working application for a decision support. A user can upload a profile for a home-loan or a personal-loan, view the results of the model, check for the contributions of SHAP and LIME, view the financial summary and keep the assessment by history or export it as a pdf file.

The most important one is the relationship between these functions. Prediction, explanation, financial arithmetic and persistence are usually studied in isolation from one another during development. It's easier to understand when they're in a single workflow as a submitted field, a displayed reason and a saved assessment value.

The distinction of the product is also crucial. The home-loan pipeline covers property context and LTV, and personal-loan pipeline includes purpose and card utilisation. These paths are coordinated by shared application logic which does not assume that their feature definitions are the same.

High performance is reflected in the test metrics saved: 98.50\% accuracy for home loans and 96.50\% accuracy for personal loans against the generated approval labels. These results are for the stratified test partitions recorded. They endorse the use of the models as part of the project's simulation, not increased likelihood of repayment.

The explanation stage is an inspectable local account of each model output. The three explanation diagnostics and separate Gender-sensitivity check reveal some aspects of that behaviour. Their assumptions remain important, and the combined Trust Score is not a universal measurement of trust.

The application also illustrates the importance of maintaining a separation between confidence and evidence for explanation. A result may have a high predicted-class probability with a further review required. Similarly, a model that has a bad dependency can be accurately described by an explanation. The final assessment should then be an output of the decision support and not an automatic authority to lend.

\section{Limitations of the Current Work}
The datasets produced include authored relationships and do not apply a Gender-dependent label adjustment. Evaluation on synthetic records will not prove generalisation and fairness for a lending population. The product evidence does not contain any external validation or probability-calibration study for repayment.

The explanation checks have limitations. Local concordance evaluates a shortened reconstruction at the original input; stability covers only the selected nearby financial changes; and consistency compares leading features and their directions. The separate fairness function compares one binary input flip. Wider validation is needed before using these checks as research-based acceptance criteria.

This current application also does not provide the server-side authentication or role enforcement. The what-if calculator can calculate affordability arithmetics without a new model prediction. Records and PDFs may be stored and reviewed, but are not a cryptographically verifiable audit trail or a lending decision signed by the lender.

\section{Future Scope}
\begin{enumerate}
\item \textbf{Evaluation with Real Lending Records:} Future work could utilize the use of representative loan records, provided with appropriate permission with personal details removed. Testing against well-defined results, such as actual repayments (where available), would help gauge the effectiveness of the system beyond more artificially created applications.

\item \textbf{Broader Model Testing:} The loan models could be compared to simpler models and tested to various applicant groups and time periods. This would assist in the detection of prediction errors, changes in the distribution of loans and whether the probabilities reported match the observed results. The generated records may further be used for tests of known failure conditions.

\item \textbf{Improved Explanation Checks:} Further evaluation might include the ability to determine whether the system provides useful explanation or misleading explanation. Adjustments during these checks should be realistic and maintain the relationships among the income, loan amounts and monthly payments. Another comparison for explanations might include the direction and magnitude of each contribution as well as differences among the groups of applicants.

\item \textbf{Trust Score Evaluation:} A future study might determine if the Trust Score is useful for bank officers to look more closely at explanations. It may be advisable to explore its weights and thresholds in the study before giving it a more powerful meaning.

\item \textbf{Secure Access and Record Management:} Future versions can have verified user accounts, access control based on the service, and more powerful privacy controls. Industry standards for retaining and discarding applicant records, combined with more reliable monitoring of storage and controlled database updating, will help to facilitate broader usage.

\item \textbf{Extended What-if Assessments:} One way in which the current affordability calculator could be expanded is by generating a new model assessment when selected application information is altered. Users could view the revised scenario and be able to compare it to the original assessment and understand exactly what changes had been made. Each scenario would be explicitly stated as a hypothetical exercise and not as an approval for a loan.
\end{enumerate}

\par\medskip\noindent
This chapter has brought together the software contribution, the available evaluation and the limits of the current system.  It has defined the next steps in data quality, model evaluation, explanations validator and application controls to provide a firmer ground for future developments.
